% Generated draft: requires mathematical and source review.
\documentclass[11pt]{article}
\usepackage[margin=25mm]{geometry}
\usepackage{fontspec}
\setmainfont{Latin Modern Roman}
\usepackage{amsmath,amssymb,mathtools,unicode-math}
\setmathfont{Latin Modern Math}
\usepackage{enumitem,needspace,xurl,hyperref}
\setlength{\parindent}{0pt}
\setlength{\parskip}{5pt}
\setlength{\emergencystretch}{3em}
\newcommand{\N}{\mathbb{N}}
\newcommand{\Z}{\mathbb{Z}}
\newcommand{\Q}{\mathbb{Q}}
\newcommand{\R}{\mathbb{R}}
\newcommand{\C}{\mathbb{C}}
\newcommand{\F}{\mathbb{F}}
\newcommand{\A}{\mathbb{A}}
\newcommand{\PP}{\mathbb{P}}
\newcommand{\E}{\mathbb{E}}
\newcommand{\eps}{\varepsilon}
\newcommand{\dd}{\,\mathrm{d}}
\begin{document}
\section*{UnsolvedMath ID: 20001954}
\textbf{Source problem code:} catalogue entry
\subsection{Definitions and mathematical statement}
The owned source is E. Lerman's Question \(3.1\) and C. Boyer's Comment \(3.2\), concerning the kernel, cokernel and possible restricted surjectivity of a contact Kirwan map. For a contact action preserving a contact form \(\alpha\), its circle moment map is \(\mu=\alpha(\xi_M)\). At a regular zero level the ordinary equivariant restriction is
\[
 \kappa:H_{S^1}^*(M;A)\longrightarrow H_{S^1}^*(\mu^{-1}(0);A),
\]
identified with quotient cohomology when the action is free. The aid uses coefficients \(A=\Z,\Q\), or \(\R\).

Lerman's example is the unit \(S^3\subset\C^2\), with standard contact form and free action
\[
 \lambda\cdot(z_1,z_2)=(\lambda z_1,\lambda^{-1}z_2),\qquad
 \mu=|z_1|^2-|z_2|^2.
\]
The full quotient is \(S^2\cong\mathbb{CP}^1\), while the zero-level quotient is \(S^1\), identified with its equator. Consequently \(\kappa\) is equator restriction
\[
 A[u]/(u^2)=H^*(S^2;A)\longrightarrow H^*(S^1;A),\qquad |u|=2,
\]
with
\[
 \ker\kappa=(u),\qquad \operatorname{coker}\kappa\cong A
       \text{ concentrated in degree }1.
\]
Cokernel here means the graded \(A\)-module cokernel, not a quotient ring by an ideal generated by the image.

The aid's broader obstruction assumes a compact connected cooriented contact manifold \(M^{2n+1}\), \(n\geq1\), a free coorientation-preserving contact circle action, an invariant contact form, and a nonempty regular zero level. Then \(B=M/S^1\) is an oriented compact \(2n\)-manifold, \(\mu\) descends to a function on \(B\), and \(Y=\mu^{-1}(0)/S^1\) is a separating oriented hypersurface. The map is \(i^*:H^*(B;A)\to H^*(Y;A)\). The pair sequence gives, degreewise,
\[
 \ker i^*=\operatorname{im}(H^k(B,Y;A)\to H^k(B;A)),
\]
\[
 \operatorname{coker}i^*\cong
 \ker(H^{k+1}(B,Y;A)\to H^{k+1}(B;A)).
\]
Orient the components of \(Y\) as the boundary of one side. Integration-sum on their top cohomology annihilates restriction from \(B\) and induces a surjection
\[
 \operatorname{coker}\bigl(H^{2n-1}(B;A)\to H^{2n-1}(Y;A)\bigr)
 \longrightarrow A.
\]
Thus this ordinary map is never fully surjective under these hypotheses. If \(Y\) is connected, restriction in that degree is zero and the cokernel is exactly \(A\). Freeness, regularity, compactness, coorientation and nonemptiness are essential parts of the reported theorem; no conclusion for all nonfree contact actions is asserted.

Boyer asks which contact or \(3\)-Sasakian classes might nevertheless admit surjectivity, possibly only up to a middle dimension as in the preceding conjecture. That earlier conjecture and its precise degree range are absent, so no numerical range is invented.

The extraction overcaptures the beginning of Section \(4\), which the supplied verification assigns to later records. For provenance only, its Goldin--Holm--Knutson theorem concerns a compact Hamiltonian torus space, regular moment value, a ring on \(\bigoplus_{g\in T}H_T^*(M^g;\Q)\), and a natural surjection to Chen--Ruan orbifold cohomology of the symplectic quotient. Its following Lerman comment distinguishes Haefliger's earlier orbifold cohomology notion. These are neighboring statements, not a contact-surjectivity theorem or part of the owned question.
\subsection{Short English statement}
For Lerman's contact three-sphere example, the map is equator restriction: its kernel is the sphere's degree-two class and its module cokernel is one degree-one copy of the coefficients. The aid extends the failure to free compact contact circle actions with nonempty regular zero quotient, using its separating-hypersurface top cohomology. Restricted and \(3\)-Sasakian variants remain questions.
\subsection{Sources}
\begin{itemize}
\item Local UnsolvedMath registry, permanent ID \(20001954\), source code \texttt{catalogue entry}; canonical collection \texttt{aim-geometry-notes.json}, notes[291].
\item AIM workshop ``Moment maps and surjectivity in various geometries,'' page \(4\), Question \(3.1\) (E. Lerman) and Comment \(3.2\) (C. Boyer): \url{https://aimath.org/WWN/momentmaps/momentmaps.pdf}. The supplied verification excludes overcaptured Section \(4\), Theorem \(4.1\), and Comment \(4.2\) from this owned item. No PDF was opened in this session.
\item Neighboring attributions preserved: Goldin--Holm--Knutson, Chen--Ruan orbifold cohomology and Lerman's Haefliger comment. No full bibliographic entries are supplied for these or the contact-map assertions. Original remarks and literature are empty, with no proposer-field value or year.
\item Supplied research attempt \(1\), \texttt{partial\_result}, run status \texttt{partially\_solved}, rigor label ``proved,'' novelty confidence ``low.'' Registry classification \texttt{PARTIAL-PROGRESS}, updated 14 August 2026.
\end{itemize}
\end{document}
